\section{Introduction}

Election-count publication is not a single computation. Election-night returns
are often unofficial; late mail ballots, provisional adjudication, ballot cures,
write-in review, duplication, recounts, and court orders may legitimately change
totals before certification. Public trust therefore depends on two ideas that
must be held together: the public should be able to replay arithmetic from
available evidence, and the software doing the replay must not pretend to make
the legal judgment that election officials and courts make.

RCOUNT is a proposed substrate for that replay. It packages election-count
records into a stable directory layout, uses canonical hashes to bind normalized
records to source evidence, and emits verification transcripts for mechanical
checks. Its target is narrower than end-to-end cryptographic voting and broader
than a single vendor export parser. RCOUNT asks: given the public artifacts an
election office is willing or required to publish, can an independent party
recompute the public count claims and locate failures at a precise layer?

The design is guided by three roles. CANVASS keeps status transitions grounded
in election administration: unofficial, canvassed, recounted, amended, and
certified totals are not interchangeable. TALLY keeps ballot and contest
accounting explicit: batches, reporting units, residual counts, write-ins, and
vendor source distinctions must survive normalization. VAULT keeps public
verification from becoming a coercion channel: RCOUNT can support inclusion and
tamper evidence, but it must not let a voter prove candidate choices.

This paper makes four contributions. First, it defines the RCOUNT package model
as a layered public-evidence artifact. Second, it describes the verifier
contract currently implemented in the Rust crates. Third, it identifies the
optional boundary with RPLAN district assignments. Fourth, it documents the
synthetic positive and negative fixture suite used to stabilize the first
contract before real state or vendor adapters are introduced.

